\documentclass[10pt]{article}
\usepackage[utf8]{inputenc}
\usepackage[vietnamese]{babel}
\usepackage{amsmath,amssymb,amsfonts}
\usepackage{geometry}
\geometry{
  paperwidth=8.5in,
  paperheight=11in,
  left=0.75in,
  right=0.75in,
  top=0.75in,
  bottom=0.75in,
  footskip=0.35in
}
\usepackage{xcolor}
\usepackage{graphicx}
\usepackage{fancyhdr}
\usepackage[most]{tcolorbox}
\usepackage{mathptmx}

% Định nghĩa màu sắc theo chuẩn sách James Stewart
\definecolor{stewartcyan}{RGB}{0, 130, 185}
\definecolor{stewartred}{RGB}{205, 32, 38}
\definecolor{stewartdarkgray}{RGB}{70, 70, 70}

% Cấu hình header & footer
\pagestyle{fancy}
\fancyhf{}
\renewcommand{\headrulewidth}{0pt}
\fancyhead[L]{%
  \fontsize{9}{11}\selectfont%
  \textbf{194}\hspace{2em}\textbf{\textsf{CHƯƠNG 3}}\hspace{1em}\textsf{Quy tắc đạo hàm}%
}

\setlength{\parindent}{0pt}
\setlength{\parskip}{0pt}

\begin{document}

% ================= KHỐI 1: BẢNG ĐẠO HÀM LƯỢNG GIÁC & GHI CHÚ CO-FUNCTIONS =================
\noindent
\begin{minipage}[t]{0.26\textwidth}
  \vspace*{28pt}
  \begin{flushleft}
    \fontsize{8.5}{11.5}\selectfont
    Khi bạn ghi nhớ bảng này, sẽ rất hữu ích nếu nhận thấy rằng dấu trừ luôn đi cùng với đạo hàm của các ``hàm co-'', nghĩa là côsin, côsec, và côtang.
  \end{flushleft}
\end{minipage}%
\hfill
\begin{minipage}[t]{0.70\textwidth}
  \fontsize{9.5}{13.5}\selectfont
  công thức đạo hàm cho các hàm lượng giác trong bảng sau. Hãy nhớ rằng chúng chỉ đúng khi $x$ được tính bằng đơn vị radian.

  \vspace{6pt}
  \begin{tcolorbox}[
    enhanced,
    colback=white,
    colframe=stewartred,
    boxrule=0.8pt,
    arc=0mm,
    left=15pt,
    right=15pt,
    top=8pt,
    bottom=10pt
  ]
    {\color{stewartred}\bfseries\textsf{\fontsize{10.5}{13}\selectfont Đạo hàm của các hàm lượng giác}}
    \vspace{0.25cm}
    
    \noindent
    \begin{minipage}[t]{0.46\textwidth}
      \begin{align*}
        \frac{d}{dx}(\sin x) &= \cos x\\[9pt]
        \frac{d}{dx}(\cos x) &= -\sin x\\[9pt]
        \frac{d}{dx}(\tan x) &= \sec^2 x
      \end{align*}
    \end{minipage}%
    \hfill
    \begin{minipage}[t]{0.50\textwidth}
      \begin{align*}
        \frac{d}{dx}(\csc x) &= -\csc x \cot x\\[9pt]
        \frac{d}{dx}(\sec x) &= \sec x \tan x\\[9pt]
        \frac{d}{dx}(\cot x) &= -\csc^2 x
      \end{align*}
    \end{minipage}
  \end{tcolorbox}
\end{minipage}

\vspace{10pt}

% ================= KHỐI 2: VÍ DỤ 2 & HÌNH 3 =================
\noindent
\begin{minipage}[t]{0.26\textwidth}
  \vspace*{72pt}
  \centering
  \includegraphics[width=\linewidth]{images/p0229_fig1.png}
  \vspace{3pt}
  \begin{flushleft}
    \fontsize{8.5}{11}\selectfont
    \textbf{\textsf{HÌNH 3}}\\ 
    Các tiếp tuyến nằm ngang trong Ví dụ 2
  \end{flushleft}
\end{minipage}%
\hfill
\begin{minipage}[t]{0.70\textwidth}
  \fontsize{9.5}{13.5}\selectfont
  \textbf{\textsf{\color{stewartcyan}VÍ DỤ 2}}\hspace{0.8em}Tính đạo hàm của $f(x) = \dfrac{\sec x}{1 + \tan x}$. Đồ thị của $f$ có tiếp tuyến nằm ngang tại các giá trị nào của $x$?

  \vspace{4pt}
  \textbf{\textsf{\color{stewartcyan}LỜI GIẢI}}\hspace{0.8em}Quy tắc thương cho ta
  \begin{align*}
    f'(x) &= \frac{(1 + \tan x)\dfrac{d}{dx}(\sec x) - \sec x \dfrac{d}{dx}(1 + \tan x)}{(1 + \tan x)^2}\\[6pt]
    &= \frac{(1 + \tan x)\sec x \tan x - \sec x \cdot \sec^2 x}{(1 + \tan x)^2}\\[6pt]
    &= \frac{\sec x(\tan x + \tan^2 x - \sec^2 x)}{(1 + \tan x)^2}\\[6pt]
    &= \frac{\sec x(\tan x - 1)}{(1 + \tan x)^2} \qquad\qquad {\color{stewartcyan}(\sec^2 x = \tan^2 x + 1)}
  \end{align*}

  \vspace{2pt}
  Vì $\sec x$ không bao giờ bằng $0$, ta thấy rằng $f'(x) = 0$ khi $\tan x = 1$, và điều này xảy ra khi $x = \pi/4 + n\pi$, với $n$ là một số nguyên (xem Hình 3).\hfill{\color{stewartcyan}$\blacksquare$}
\end{minipage}

\vspace{10pt}

% ================= KHỐI 3: VÍ DỤ 3 & HÌNH 4 =================
\noindent
\begin{minipage}[t]{0.26\textwidth}
  \vspace*{78pt}
  \centering
  \includegraphics[width=0.82\linewidth]{images/p0229_fig2.png}
  \vspace{4pt}
  \begin{flushleft}
    \fontsize{8.5}{11}\selectfont
    \textbf{\textsf{HÌNH 4}}
  \end{flushleft}
\end{minipage}%
\hfill
\begin{minipage}[t]{0.70\textwidth}
  \fontsize{9.5}{13.5}\selectfont
  Các hàm lượng giác thường được sử dụng để mô hình hóa các hiện tượng thực tế. Đặc biệt, các dao động, sóng, chuyển động đàn hồi và các đại lượng khác biến thiên có tính tuần hoàn đều có thể được mô tả bằng cách sử dụng các hàm lượng giác. Trong ví dụ tiếp theo, chúng ta sẽ khảo sát một trường hợp dao động điều hòa đơn giản.

  \vspace{7pt}
  \textbf{\textsf{\color{stewartcyan}VÍ DỤ 3}}\hspace{0.8em}Một vật gắn vào đầu của một lò xo thẳng đứng được kéo dãn xuống dưới $4\text{ cm}$ khỏi vị trí cân bằng rồi thả ra tại thời điểm $t = 0$. (Xem Hình 4 và chú ý rằng chiều hướng xuống là chiều dương.) Vị trí của vật tại thời điểm $t$ là
  \[
    s = f(t) = 4\cos t
  \]
  Hãy tìm vận tốc và gia tốc tại thời điểm $t$, sau đó sử dụng chúng để phân tích chuyển động của vật.

  \vspace{4pt}
  \textbf{\textsf{\color{stewartcyan}LỜI GIẢI}}\hspace{0.8em}Vận tốc và gia tốc lần lượt là
  \begin{align*}
    v &= \frac{ds}{dt} = \frac{d}{dt}(4\cos t) = 4\frac{d}{dt}(\cos t) = -4\sin t\\[6pt]
    a &= \frac{dv}{dt} = \frac{d}{dt}(-4\sin t) = -4\frac{d}{dt}(\sin t) = -4\cos t
  \end{align*}
\end{minipage}

\vfill
\noindent
\begin{center}
\fontsize{6.8}{8.8}\selectfont\color{stewartdarkgray}
Copyright 2021 Cengage Learning. All Rights Reserved. May not be copied, scanned, or duplicated, in whole or in part. WCN 02-200-203\\ 
Due to electronic rights, some third party content may be suppressed from the eBook and/or eChapter(s). Editorial review has deemed that any suppressed content does not materially affect the overall learning experience. Cengage Learning reserves the right to remove additional content at any time if subsequent rights restrictions require it.
\end{center}

\end{document}